\documentclass[11pt]{article}
\usepackage[margin=1in]{geometry}

% Core packages
\usepackage{amsmath,amssymb}
% \usepackage{tikz-cd}
\usepackage{multicol}
\usepackage{array}
\usepackage{booktabs}
\usepackage{longtable}
\usepackage{enumitem}
\usepackage{graphicx} % To include figures (Gantt, Pie chart, etc.)
\usepackage[hidelinks]{hyperref}

\renewcommand{\tablename}{Tabla}
\renewcommand{\contentsname}{Índice}
\renewcommand{\refname}{Bibliografía}
\renewcommand{\figurename}{Figura}

% Paragraphs
\setlength{\parindent}{0pt}
\setlength{\parskip}{1\baselineskip}

\title{Informe del Proyecto Programado\\Biblioteca UCR Recinto de Grecia}
\author{Marcos Ferreto Estrada \and Paulo Anchía Correás}
\date{\today}

\begin{document}

\begin{titlepage}
    \centering
    {\Large Universidad de Costa Rica\par}
    {\large Sede de Occidente - Recinto de Grecia\par}
    {\large Curso: Estructuras de Datos\par}

    \vspace{2.5cm}
    {\LARGE\bfseries Informe del Proyecto Programado\par}
    {\Large Sistema de administración de libros de biblioteca\par}

    \vspace{2.5cm}
    \begin{tabular}{rl}
        Integrantes: & Marcos Ferreto Estrada --- C5F092 \\
                     & Paulo Anchía Correás --- [Escribir Carné] \\
        Profesora:   & [Nombre de la profesora] \\
        Grupo:       & [Número de grupo] \\
        Fecha:       & \today \\
    \end{tabular}

    \vfill
\end{titlepage}

\tableofcontents
\newpage

\section{Introducción}

[Escribir aquí una introducción. Describir brevemente el propósito del documento y del proyecto de biblioteca.]

\section{Descripción del problema}

[Explicar aquí con sus propias palabras qué hace el sistema: administrar libros, estudiantes y préstamos usando estructuras de datos. Mencionar que permite buscar, insertar, eliminar y registrar préstamos, con persistencia y GUI.]

\section{Análisis del problema}

\subsection{Estructuras de datos utilizadas}

\subsubsection{Árbol AVL para libros}

[Explicar aquí por qué se usó un Árbol AVL para los libros. Ej: porque permite búsquedas rápidas y mantiene el balance automáticamente.]

\textbf{Datos almacenados por libro:}
\begin{itemize}[noitemsep]
    \item [Escribir dato 1]
    \item [Escribir dato 2]
\end{itemize}

\subsubsection{Tabla hash para estudiantes}

[Explicar aquí por qué se usó una Tabla Hash para los estudiantes. Ej: facilita búsquedas por carnet y reduce el tiempo de acceso.]

\textbf{Datos almacenados por estudiante:}
\begin{itemize}[noitemsep]
    \item [Escribir dato 1]
    \item [Escribir dato 2]
\end{itemize}

\subsubsection{Árbol rojinegro para préstamos}

[Explicar aquí por qué se usó un Árbol Rojinegro para los préstamos. Ej: maneja bien inserciones y consultas sin degradarse.]

\textbf{Datos almacenados por préstamo:}
\begin{itemize}[noitemsep]
    \item [Escribir dato 1]
    \item [Escribir dato 2]
\end{itemize}

\subsection{Carga y persistencia de archivos}

[Describir aquí cómo el programa lee y guarda los datos en los archivos de texto/XML. Mencionar que se actualizan al modificar registros.]

\subsection{Resolución de subproblemas}

[Completar la tabla explicando cómo resolvieron cada parte del proyecto.]

\begin{longtable}{p{0.35\textwidth} p{0.55\textwidth}}
\toprule
\textbf{Subproblema} & \textbf{Solución implementada} \\
\midrule
Buscar libros por autor & [Explicar solución] \\
Buscar libro por título & [Explicar solución] \\
Buscar libro por código & [Explicar solución] \\
Eliminar libro por código & [Explicar solución] \\
Buscar estudiante por carné & [Explicar solución] \\
Buscar estudiante por nombre & [Explicar solución] \\
Buscar estudiantes por carrera & [Explicar solución] \\
Eliminar estudiante por carné & [Explicar solución] \\
Realizar préstamo & [Explicar solución] \\
Mostrar árboles en inorden & [Explicar solución] \\
Interfaz gráfica & [Explicar solución] \\
\bottomrule
\end{longtable}

\subsection{Funciones o métodos principales}

[Incluir explicación breve de las funciones más relevantes del código.]

\begin{longtable}{p{0.25\textwidth} p{0.24\textwidth} p{0.18\textwidth} p{0.23\textwidth}}
\toprule
\textbf{Función o método} & \textbf{Parámetros} & \textbf{Retorna} & \textbf{Descripción} \\
\midrule
\texttt{cargar\_todo()} & [Parámetros] & [Retorno] & [Descripción breve] \\
\texttt{guardar\_todo()} & [Parámetros] & [Retorno] & [Descripción breve] \\
\texttt{insertar\_libro()} & [Parámetros] & [Retorno] & [Descripción breve] \\
\texttt{buscar\_por\_carnet()} & [Parámetros] & [Retorno] & [Descripción breve] \\
\texttt{crear\_prestamo()} & [Parámetros] & [Retorno] & [Descripción breve] \\
\texttt{eliminar\_libro()} & [Parámetros] & [Retorno] & [Descripción breve] \\
\texttt{eliminar\_estudiante()} & [Parámetros] & [Retorno] & [Descripción breve] \\
\bottomrule
\end{longtable}

\section{Análisis de resultados}

[Completar la tabla con el avance del proyecto y observaciones.]

\begin{longtable}{p{0.38\textwidth} p{0.24\textwidth} p{0.28\textwidth}}
\toprule
\textbf{Parte del proyecto} & \textbf{Estado} & \textbf{Observaciones} \\
\midrule
Árbol AVL & Completado & [Observación] \\
Tabla hash & Completado & [Observación] \\
Árbol rojinegro & Completado & [Observación] \\
Préstamos & Completado & [Observación] \\
Eliminación segura & Completado & [Observación] \\
Interfaz gráfica de usuario & Completado & [Observación] \\
Informe escrito & En progreso & [Observación] \\
\bottomrule
\end{longtable}

\section{Conclusiones y recomendaciones}

[Escribir al menos 6 conclusiones y sus respectivas recomendaciones. A continuación, un formato base para llenarlas.]

\begin{enumerate}
    \item \textbf{Conclusión:} [Escriba la primera conclusión. Ej: Las estructuras de datos mejoran la eficiencia del sistema.]\\
    \textbf{Recomendación:} [Escriba la recomendación asociada.]

    \item \textbf{Conclusión:} [Escriba la segunda conclusión. Ej: El AVL fue útil para mantener orden y rapidez en libros.]\\
    \textbf{Recomendación:} [Escriba la recomendación asociada.]

    \item \textbf{Conclusión:} [Escriba la tercera conclusión. Ej: La tabla hash facilitó el acceso a estudiantes por carnet.]\\
    \textbf{Recomendación:} [Escriba la recomendación asociada.]

    \item \textbf{Conclusión:} [Escriba la cuarta conclusión. Ej: El árbol rojinegro permitió manejar préstamos de forma estable.]\\
    \textbf{Recomendación:} [Escriba la recomendación asociada.]

    \item \textbf{Conclusión:} [Escriba la quinta conclusión. Ej: La validación antes de eliminar evitó inconsistencias.]\\
    \textbf{Recomendación:} [Escriba la recomendación asociada.]

    \item \textbf{Conclusión:} [Escriba la sexta conclusión. Ej: La GUI hizo el sistema más usable para el usuario final.]\\
    \textbf{Recomendación:} [Escriba la recomendación asociada.]
\end{enumerate}

\section{Bibliografía}

[Incluyan las fuentes de estructuras de datos y documentación usada en formato APA.]

\begin{itemize}
    \item [Escribir fuente 1. Ej: Python Software Foundation. (2026). tkinter...]
    \item [Escribir fuente 2. Ej: Cormen, T. H. ...]
    \item [Escribir fuente 3]
\end{itemize}

\section{Anexos (Diagramas)}

[Insertar aquí el código o imágenes para los diagramas requeridos:]
\begin{itemize}
    \item Diagrama de Gantt (planificación del trabajo).
    \item Pie chart (distribución del esfuerzo por tareas).
    \item Workload (carga por integrante).
    \item Burndown (avance del proyecto por fechas).
\end{itemize}

\end{document}
